{\scriptsize
        \begin{tabularx}{\textwidth}{|p{0.25\columnwidth}|X|}
            \hline 
            \textbf{Notions et contenus} & \textbf{Capacités exigibles} \\
            \hline 
            \multicolumn{2}{|l|}{\textbf{6.1.2. Ondes acoustiques dans les fluides }} \\
            \hline
            Approximation acoustique. Équation de d'Alembert pour la surpression acoustique.  & 
            Classer les ondes acoustiques par domaines fréquentiels.
Valider l'approximation acoustique. Établir, par une approche eulérienne, l'équation de propagation de la surpression acoustique dans une situation unidimensionnelle en coordonnées cartésiennes.
Utiliser l'opérateur laplacien pour généraliser l'équation d'onde. \\
            \hline 
             Célérité des ondes acoustiques.   & 
            Exprimer la célérité des ondes acoustiques en fonction de la température pour un gaz parfait. \\
        \hline 
        Ondes planes progressives harmoniques : caractère longitudinal, impédance acoustique. & 
        Exploiter la notion d'impédance acoustique pour faire le lien entre les champs de surpression et de vitesse d'une onde plane progressive harmonique.
Utiliser le principe de superposition des ondes planes progressives harmoniques. \\ 
        \hline
        Densité volumique d'énergie acoustique, vecteur densité de courant énergétique. Intensité sonore. Niveau d'intensité sonore.  & 
        Utiliser les expressions admises du vecteur densité de courant énergétique et de la densité volumique d'énergie associés à la propagation de l'onde.
        Citer quelques ordres de grandeur de niveaux d'intensité sonore. \\
        \hline 
        Ondes acoustiques sphériques harmoniques.   & 
        Utiliser une expression fournie de la surpression pour interpréter par un argument énergétique la décroissance en $1 / r$ de l'amplitude. \\
        \hline
        \end{tabularx}
\vspace{0.5cm}
}
